% =====================================================================
% EJE B - LAS TECNICAS DEL TRACK A  ---  RESPONSABLE: INTEGRANTE 2
% Aqui se justifica el track CON LITERATURA, no con opinion.
% De cada familia: fortalezas, limites y SUPUESTOS bajo los que funciona.
% =====================================================================
\subsection{Axis B: Techniques}
\label{sec:ejeB}

\subsubsection*{\textbf{Gradient Boosting and Random Forests versus neural networks on tabular
data}} Tree-based ensembles are strong candidates for the present project because
they can capture nonlinear relationships and interactions between variables
without requiring a predefined functional form. The literature also provides
evidence that their performance on structured tabular data can be competitive
with, and in some settings superior to, more complex models. Hall and
Rasheed~\cite{ejeB_hall2025} conduct a systematic review of 79 studies comparing
tree-based machine learning (TBML), including Random Forest, XGBoost, LightGBM,
and CatBoost, with deep learning across multiple domains. TBML methods are the
most frequent top performers in their review and require substantially less
training time on average, with reported differences of two to four orders of
magnitude. However, their evidence comes primarily from time-series forecasting
rather than horse-racing prediction, and the authors emphasize that performance
depends on factors such as dataset characteristics, temporal structure, and
experimental design. Therefore, the review supports considering tree-based
methods as strong baseline candidates, but does not establish their superiority
for the present dataset.

Evidence from other structured-data domains further supports this choice while
illustrating the importance of empirical comparison. Dumitrescu et
al.~\cite{ejeB_dumitrescu2022} compare Random Forest, logistic regression, SVM,
and neural networks across four credit-scoring datasets. On the Kaggle
\textit{Give Me Some Credit} dataset, for example, Random Forest and their
tree-augmented logistic model PLTR achieve AUC values of 0.853 and 0.857,
respectively, compared with 0.752 for the neural network. These results show
that a more complex model does not necessarily provide better predictive
performance on structured data. Similarly, Yuan et al.~\cite{ejeB_yuan2025}
review machine-learning approaches for sports-injury prediction and report
strong results from gradient-boosted tree methods while also identifying
settings in which simpler logistic-regression models perform comparably or
better. Together, these studies support treating model complexity as an
experimental variable rather than assuming that neural networks will provide
the best performance. Within racing-related applications, Kholkine et
al.~\cite{ejeB_kholkine2020} use XGBoost to predict road-cycling race
performance. The authors motivate the use of XGBoost by its predictive
performance, computational efficiency, and ability to handle missing values.
Their application is particularly relevant to the present study because
missing values can represent meaningful events, such as riders who did not
finish a previous race. This illustrates why missing historical information
should be analyzed according to its meaning before an imputation strategy is
selected.

\subsubsection*{\textbf{Ranking within a group, not independent binary classification}} A horse
race produces a total order over a small, closed set of competitors, which is a
ranking problem, not a set of independent win/lose labels, yet none of the
horse-racing-specific studies we read in \ref{sec:ejeA} model it that way;
all of them predict a per-horse label or score and then sort. This distinction is
important when defining the prediction target: a model that estimates the
probability of winning for each horse is different from a model whose objective
is to predict the relative order of all horses within a race.

The distinction between pointwise prediction and group-level ranking is illustrated
by work outside and within horse racing. Peng et al.~\cite{ejeB_peng2020} formulate
rank-position forecasting in IndyCar as the prediction of future rank positions
within a race. Their baselines include pointwise machine-learning regression
models such as Random Forest, SVM, and XGBoost, as well as a naive baseline
called CurRank, which assumes that the current rank will not change. Their proposed
RankNet explicitly models rank-position sequences and reports improvements over the
baselines in the evaluated forecasting tasks. In particular, the authors report
that RankNet-MLP exceeds the CurRank baseline by 7\% in Top-1 accuracy and 19\% in
MAE for the main forecasting task. These results do not demonstrate that ranking
models are necessarily superior for horse racing, but they show that the formulation
of the target and the structure of the competitors can materially affect the
prediction problem.

Within horse racing, Silverman and Suchard~\cite{ejeA_gupta2023} provide a
complementary approach through regularized conditional logistic regression with
frailty, which incorporates dependence among competitors within the same race when
predicting the winner. Although their model does not predict the complete finishing
order, it demonstrates that race-level structure can be incorporated directly into
the statistical formulation rather than treating each horse as an entirely
independent observation.

These studies suggest that the race should be considered the natural grouping unit
when defining the prediction task and evaluating the resulting model. However, they
do not by themselves establish which formulation is most appropriate for the
present dataset. A binary classification formulation is suitable if the objective
is to estimate whether an individual horse wins its race. If the objective is to
produce an ordering of the competitors, a ranking-oriented formulation would be
more directly aligned with the target.

\subsubsection*{\textbf{Calibration and explainability}} If the final system produces probabilities
rather than only a predicted ordering, discrimination metrics alone are
insufficient. Vuko et al.~\cite{ejeB_vuko2025} explicitly evaluate calibration in
addition to discrimination and argue that model selection should consider whether
predicted probabilities correspond adequately to observed frequencies. Accordingly,
log loss, Brier score, or an appropriate calibration analysis can complement
ranking and classification metrics in the experimental protocol.

For interpretability, Kholkine et al.~\cite{ejeB_kholkine2020} use SHAP-based
feature attribution to inspect the contribution of predictors in their XGBoost
model. SHAP should be regarded as an interpretability tool rather than as evidence
that XGBoost is inherently superior. Its role in project is to analyze which
pre-race features contribute to model predictions after a model has been validated,
while avoiding causal interpretations from feature-attribution values alone.

\subsubsection*{\textbf{Class imbalance and evaluation}} Predicting the winner of a race produces
an intrinsically imbalanced target when each horse is represented as an individual
observation. Gupta and Singh~\cite{ejeA_gupta2023} illustrate this problem clearly:
only 11.8\% of their observations correspond to winners. After applying SMOTE,
tetthey report substantially higher F1 and ROC values for Random Forest than in the
unbalanced setting. However, synthetic oversampling must be treated carefully when
observations belong to races, because generating synthetichorses without respecting
race membership can distort the competitive structure. Hence, imbalance-handling
methods such as class weighting or resampling should be evaluated experimentally
rather than adopted solely because they improved a result in another dataset.

Overall, the researches supports a candidate set centered on logistic regression
and tree ensembles, with Random Forest and gradient boosting methods such as XGBoost
as flexible nonlinear alternatives. KNN, Naive Bayes, AdaBoost, and Bagging can be
used as secondary comparison models if computational and experimental scope allow.
The final selection should be determined empirically using a common temporal
protocol and metrics aligned with the final target definition. This approach avoids
assuming that the most complex or most frequently reported algorithm is necessarily
the best choice for the Hong Kong racing dataset.

For these reasons, the selected track is classical machine learning over structured
tabular data. Consequently, the relevant question is not which algorithm is
universally the best, but which model families provide an appropriate balance
between predictive capacity, assumptions, computational cost, and interpretability
for the structure of horse-racing data. The literature reviewed for this project
suggests that tree ensembles, regularized linear models, and other classical
classifiers are viable candidates, while model choice should remain dependent on the
target definition, feature representation, and evaluation protocol.